% Source: https://artofproblemsolving.com/community/c3h2047101_domain_and_range_of_a_rational_function
% Stable identity: aops-topic-2047101-q01
Both denominators are at least \(1\), so both functions are defined for every real \(x\).

\textbf{(a)} A number \(y\) is in the range exactly when \(y(x^2+1)=x\), i.e. \(yx^2-x+y=0\), has a real solution \(x\). If \(y=0\), then \(x=0\) works. If \(y\ne0\), this quadratic has a real root iff its discriminant is nonnegative:
\[
1-4y^2\ge 0\iff -\dfrac{1}{2}\le y\le \dfrac{1}{2} .
\]
So the range is \(\boxed{\left[-\dfrac{1}{2},\dfrac{1}{2}\right]}\), attained at \(x=-1\) and \(x=1\).

\textbf{(b)} \(h\) is continuous on \(\mathbb{R}\). Since \(h(x)=x-\dfrac{x}{x^2+1}\) and \(\left|\dfrac{x}{x^2+1}\right|\le\dfrac{1}{2}\), we have \(h(x)\ge x-\dfrac{1}{2}\) and \(h(x)\le x+\dfrac{1}{2}\). So \(h(x)\to+\infty\) as \(x\to+\infty\) and \(h(x)\to-\infty\) as \(x\to-\infty\). Given any real \(y\), \(h(y+1)\ge y+\dfrac{1}{2}>y\) and \(h(y-1)\le y-\dfrac{1}{2}<y\), so by the intermediate value theorem \(h(x)=y\) for some \(x\) between \(y-1\) and \(y+1\). Hence the range of \(h\) is \(\boxed{\mathbb{R}}\).

So a rational function defined on all of \(\mathbb{R}\) may have a restricted range (part (a)) or the whole real line (part (b)).
